% ============================================================
% Face-Aligned Reading-C Invariant-Subspace Structural Theorem
%   for the Net DI-Bit Current at the Host Vertex
%   under (H1_F) Face-Aligned Reading C in the 600-Cell
% Publication-Grade Hardening (Symmetry-Only Structural Closure)
%   of the OPEN-FP-F1-5 Trajectory's Sequence-1B Sub-Question
% Conscious Point Physics Flagship Paper Series
%   (F.1 Dynamical Substrate Law sub-question hardened-theorem artifact;
%    THIRTEENTH artifact in the F.1 Layer 3 hardened-theorem sequence —
%    second Sequence-1 artifact under the OPEN-FP-F1-5 non-vertex-aligned
%    Reading-C variant trajectory; closes the structural sub-question of
%    OPEN-FP-F1-5 Sequence-1B at publication-grade Layer 3 unconditional.
%    The result is a 3-dimensional invariant-subspace conclusion —
%    substantively the WEAKEST of the three Reading-C variants and
%    substantively distinct from both the vertex-aligned 1D result
%    (THEO-DSL-4 / Patch 0585) and the edge-aligned 2D result
%    (THEO-DSL-6 / Patch 0596).)
%
% Patch 0597 (Session 146, 27 May 2026) -- second OPEN-FP-F1-5 trajectory
%   closure Patch under the session-open decision-gate bundle DG-F15-1
%   through DG-F15-5 accepted as-is by Thomas at Session 146; completes
%   DG-F15-1 default-(C) priority of front-loading both Sequence-1
%   structural findings before either Sequence-2 coefficient enumeration.
% ============================================================
% Version 1.2 --- 17 August 2026 (Patch 3214):
%   added a How-we-review methodology note stating plainly that AI panel
%   review is a self-applied verification method and not journal peer
%   review; twelve chirality papers retitled so the descriptive title
%   leads and the identifier is demoted to a subtitle. No claim,
%   equation, number or section changed.
% Version 1.1 --- 17 August 2026 (Patch 3213):
%   extended the generated identifier appendix to all five identifier
%   families (THEO-, FI-, PRED-, CONJ-, PH- as well as OPEN-), and
%   replaced review-convergence scores with AI review panel wording
%   where present. No claim, equation, number or section changed.
% Version 1.0 --- 17 August 2026 (Patch 3212):
%   added the generated plain-language appendix glossing the internal
%   programme identifiers used in this paper (founder jargon ruling, 16
%   Aug 2026). No claim, equation, number or section changed.
%   This is the FIRST version stamp assigned to this paper; 1.0 denotes
%   the first stamped revision, NOT a claim of v1.0 shipped grade.

\documentclass[11pt,letterpaper]{article}
\usepackage[utf8]{inputenc}
\usepackage[margin=1in]{geometry}
\usepackage[T1]{fontenc}
\usepackage{amsmath,amssymb,amsfonts,amsthm}
\usepackage{xcolor}
\usepackage{hyperref}
\usepackage{enumitem}
\usepackage{parskip}
\usepackage{mdframed}

\hypersetup{
    colorlinks=true,
    linkcolor=blue,
    citecolor=blue,
    urlcolor=blue
}

\newtheorem{theorem}{Theorem}[section]
\newtheorem{proposition}[theorem]{Proposition}
\newtheorem{lemma}[theorem]{Lemma}
\newtheorem{corollary}[theorem]{Corollary}
\newtheorem{definition}[theorem]{Definition}
\newtheorem{remark}[theorem]{Remark}

\newcommand{\nhat}{\hat{n}}
\newcommand{\nhatperp}{\hat{n}_{\text{F}\perp}}
\newcommand{\nhatmid}{\hat{n}_{\text{ax}}}
\newcommand{\nhatrho}{\hat{n}_{\rho}}
\newcommand{\vhost}{v_{\text{host}}}
\newcommand{\ui}{u_i}
\newcommand{\uj}{u_j}
\newcommand{\Gsixcell}{\Gamma_{600}}
\newcommand{\Reals}{\mathbb{R}}
\newcommand{\Hthree}{H_3}
\newcommand{\Hfour}{H_4}
\newcommand{\Icoh}{I_h}
\newcommand{\Cs}{C_s}
\newcommand{\Dfive}{D_5}
\newcommand{\Aone}{A_1}
\newcommand{\Aprime}{A'}
\newcommand{\Aprpr}{A''}
\newcommand{\Oof}[1]{\mathcal{O}(#1)}
\newcommand{\jDI}{\vec{j}_{DI}^{\,\text{net}}}
\newcommand{\Ek}{E_k}

\title{Face-Aligned Reading-C Invariant-Subspace Structural Theorem \\
for the Net DI-Bit Current at the Host Vertex \\
in the 600-Cell \\
\large Symmetry-Only Publication-Grade Layer~3 Closure \\
of the Sequence-1B Structural Sub-Question of OPEN-FP-F1-5}

\author{Thomas Lee Abshier, ND, and Claude Opus 4.7 \\ Hyperphysics Institute --- F.1 Dynamical Substrate Law trajectory}

\date{27 May 2026 (Session 146, CPP Patch 0597)\\[2pt]
{\small Version~1.0 --- 17 August 2026 (Patch 3212: added the generated plain-language appendix glossing the internal programme identifiers used in this paper (founder jargon ruling, 16 Aug 2026). No claim, equation, number or section changed.)}\\[2pt]
{\small Version~1.1 --- 17 August 2026 (Patch 3213: extended the generated identifier appendix to all five identifier families (THEO-, FI-, PRED-, CONJ-, PH- as well as OPEN-), and replaced review-convergence scores with AI review panel wording where present. No claim, equation, number or section changed.)}\\[2pt]
{\small Version~1.2 --- 17 August 2026 (Patch 3214: added a How-we-review methodology note stating plainly that AI panel review is a self-applied verification method and not journal peer review; twelve chirality papers retitled so the descriptive title leads and the identifier is demoted to a subtitle. No claim, equation, number or section changed.)}}

\begin{document}
\maketitle

\begin{abstract}
We give a publication-grade Layer~3 \emph{symmetry-only} closure of the structural sub-question of OPEN-FP-F1-5 Sequence~1B: under (H1$_F$)~face-aligned Reading~C with $\nhat = \nhatperp$ chosen perpendicular to a triangular face $\{\vhost, \ui, \uj\}$ incident to the host vertex $\vhost$, (H2)~600-cell unit-vertex normalisation, (H3$_F$)~the residual symmetry $\Cs$ at $\vhost$ fixing the chosen face and $\nhatperp$ direction, (H4)~the F.1 v1.0 Theorem~6.1 perturbation-locality propagation rule, and (H5$_F$)~the Mechanism~A perturbation $r(\hat{e}) = r_0(1 + \delta\,\hat{e}\cdot\nhatperp)$ on each oriented edge of the 600-cell edge graph at first order in $\delta$ with path-integral extension at $\Oof{\delta^k}$ at sketch level per DG-F15-4 default~(A), the net DI-bit current at the host vertex at every non-zero order $k \geq 1$ in $\delta$ lies in the \emph{three-dimensional} $\Cs$-invariant subspace of $\Reals^4$ at $\vhost$:
\begin{equation*}
\jDI(\vhost)\big|_{\Oof{\delta^k}} \in \operatorname{span}\bigl\{\nhatrho,\,\nhatmid,\,\nhatperp\bigr\} \qquad \text{for all } k \geq 1,
\end{equation*}
where $\nhatrho := \vhost/|\vhost|$ is the radial direction at $\vhost$, $\nhatmid$ is the in-face axis from $\vhost$ to the midpoint of the chosen face edge $\{\ui, \uj\}$, and $\nhatperp$ is the face-perpendicular Reading-C direction. Equivalently, $\jDI(\vhost)|_{\Oof{\delta^k}}$ is constrained only to be perpendicular to the single $\Cs$-antisymmetric direction $\ui - \uj$.

This is the \emph{weakest of the three Reading-C variants}: the structural constraint reduces from 4-vector to 3-vector (loss of one degree of freedom), compared to the edge-aligned 2-vector (loss of two; Patch~0596 / THEO-DSL-6) and the vertex-aligned 1-vector (loss of three; Patch~0585 / THEO-DSL-4). The dimensional-growth pattern $1 \to 2 \to 3$ as the stabiliser shrinks $\Icoh \to \Dfive \to \Cs$ is the central physical insight of OPEN-FP-F1-5: \textbf{the strong parallel-to-$\nhat$ substrate-locality of F.1 Theorem~7.1 is a vertex-aligned-Reading-C-specific consequence of the maximal residual symmetry $\Icoh$, with monotonic structural weakening as the residual symmetry decreases}.

The proof factors into Lemma~\ref{lem:cs_invariance} ($\Cs$-invariance of $\jDI(\vhost)$ order-by-order via $\Cs$-equivariance of Mechanism~A's rate function under stabilisation of $\nhatperp$) and Lemma~\ref{lem:cs_invariant_subspace} (the $\Cs$-invariant subspace of $\Reals^4$ at $\vhost$ is the 3-dimensional hyperplane perpendicular to $\ui - \uj$, via the direct decomposition $\Reals^4 = \Aprime^{\oplus 3} \oplus \Aprpr^{\oplus 1}$ under $\Cs$). This artifact is the \textbf{thirteenth} of the F.1 Layer~3 hardened-theorem sequence (after Patches 0550 + 0551 + 0552 + 0571 + 0585 + 0587 + 0588 + 0589 + 0590 + 0591 + 0592 + 0596), and the \textbf{second} under the OPEN-FP-F1-5 trajectory (closing Sequence~1B after Sequence~1A's edge-aligned closure at Patch~0596). Per DG-F15-5 default~(A) the result is registered as candidate \textbf{THEO-DSL-8} entry (THEO-DSL-7 is reserved for the Sequence-2A candidate-coefficient closure under (H5$_E$)). An \textbf{anti-erasure remark} (\S\ref{sec:antierasure}) registers a stabiliser-identification discrepancy with \texttt{frontier\_sectors/FP.md} line~230 (which states ``residual $D_2$ symmetry at face centroid'') -- the correct stabiliser under the vertex-host convention (preserved from Sequence-1A precedent) is $\Cs$ of order~2, not $D_2$ of order~4; under the alternative face-centroid-host convention the stabiliser would be $D_3$ of order~6.
\end{abstract}

\section{Introduction and statement of the main result}\label{sec:introduction}

\subsection{Background and continuity with Sequence~1A}\label{subsec:background}

The OPEN-FP-F1-5 trajectory opened at Patch~0595 with a scoping document identifying three Reading-C variants (vertex-aligned, edge-aligned, face-aligned) and proposing a Sequence-1 (symmetry-only structural) + Sequence-2 (closed-form coefficient) trajectory structure for each non-vertex-aligned variant. The Session~146 open decision-gate bundle DG-F15-1 default~(C) accepted by Thomas requires front-loading both Sequence-1 structural findings before either Sequence-2 coefficient enumeration. Sequence-1A (edge-aligned) closed at Patch~0596 with a 2-dimensional invariant subspace finding under $\Dfive$ stabiliser. The present Patch~0597 closes Sequence-1B (face-aligned), completing DG-F15-1 (C)'s priority of having both Sequence-1 structural findings in hand before proceeding to Sequence-2 coefficient work.

The face-aligned variant is the most aggressive symmetry-breaking step from vertex-aligned among the three Reading-C variants, and the structural finding is correspondingly the weakest constraint on $\jDI(\vhost)$. Whether this conclusively establishes vertex-aligned Reading C as the uniquely structurally-strong variant is the substantive question addressed by Theorem~\ref{thm:main} below.

\subsection{Setup and host-convention selection}\label{subsec:setup}

Per Sequence-1A precedent (Patch~0596 \S{}1.2), we preserve the vertex-host convention: $\vhost$ remains a 600-cell vertex (not a face-centroid). The Reading-C primitive $\nhat$ varies among the three Reading-C choices only in its \emph{direction} at $\vhost$, while the host point itself remains a 600-cell vertex throughout the F.1 substrate-locality framework. This is consistent with F.1 v1.0's vertex-host scope and avoids cross-host complications from the moved-host alternative.

\textbf{Choice of face.} Choose a triangular face $F = \{\vhost, \ui, \uj\}$ of the 600-cell incident to $\vhost$, where $\ui$ and $\uj$ are two adjacent first-shell neighbours of $\vhost$ that are themselves adjacent to each other (forming the face). The 600-cell has $1200/120 \cdot 3 / 1 = 30$ triangular faces incident to each vertex, by orbit-stabiliser; the choice of face is arbitrary and the structural conclusion is face-independent by transitivity of $\Icoh$ on faces incident to $\vhost$.

\textbf{Face-perpendicular direction.} The face-aligned Reading-C primitive is
\begin{equation}\label{eq:nperp}
\nhatperp \;:=\; \frac{F_\perp}{|F_\perp|}, \qquad F_\perp \text{ a unit vector in } \Reals^4 \text{ perpendicular to the 2D affine face plane through } F.
\end{equation}
The 2D affine face plane is determined by $\{\vhost, \ui, \uj\}$; its perpendicular complement in $\Reals^4$ is 2D. By the 600-cell's symmetry, the natural choice of $\nhatperp$ is the direction from the polytope centre to the face centroid; equivalently, $\nhatperp \propto (\vhost + \ui + \uj)/3$ normalised. Under this convention $\nhatperp$ is orthogonal to the in-face direction $\ui - \uj$ and to the in-face axis $\nhatmid$ defined below.

\textbf{Radial direction at host.} As in Sequence-1A:
\begin{equation}\label{eq:nrho}
\nhatrho \;:=\; \vhost/|\vhost| \;=\; \vhost \quad \text{(under (H2) unit-vertex normalisation)}.
\end{equation}

\textbf{In-face axis at host.} The third distinguished direction at $\vhost$ relevant to the face-aligned setup is the in-face direction from $\vhost$ toward the midpoint $(\ui + \uj)/2$ of the face's non-host edge:
\begin{equation}\label{eq:nmid}
\nhatmid \;:=\; \frac{(\ui + \uj) - 2\vhost}{|(\ui + \uj) - 2\vhost|}.
\end{equation}
This is a unit vector tangent to $S^3$ at $\vhost$, in the face plane, and perpendicular to $\ui - \uj$ (the within-face edge connecting the two non-host face vertices).

\subsection{Stabiliser identification}\label{subsec:stabiliser}

\textbf{Stabiliser of $\vhost$ + chosen face.} The stabiliser in $\Icoh$ of the unordered face $F = \{\vhost, \ui, \uj\}$ (i.e., the subgroup of $\Icoh = \text{stab}_{\Hfour}(\vhost)$ that maps $F$ to itself as a set) is computed via orbit-stabiliser. With $|\Icoh| = 120$ and 30~faces incident to $\vhost$, the unordered-face stabiliser has order $120/30 = 4$. This is a $V_4$ (Klein four-group) structure: identity, $C_2$ rotation around the in-face axis from $\vhost$ to edge-midpoint, reflection in the plane containing $\vhost$ and the radial direction perpendicular to $\ui - \uj$, and reflection in the face plane itself.

\textbf{Subgroup preserving oriented $\nhatperp$.} Among the four $V_4$ elements, two flip $\nhatperp \to -\nhatperp$ (those whose action on the face's 2D orthogonal complement includes a flip) and two preserve $\nhatperp$ direction. The two $\nhatperp$-preserving elements are:
\begin{itemize}[leftmargin=*]
    \item Identity $e$;
    \item The reflection $\sigma_E$ in the 3-hyperplane perpendicular to $\ui - \uj$ through the midpoint of $\{\ui, \uj\}$. This reflection swaps $\ui \leftrightarrow \uj$, fixes $\vhost$ (because $\vhost$ is equidistant from $\ui$ and $\uj$), and preserves $\nhatperp$ direction (because the reflection's fixed 3-hyperplane contains $\nhatperp$).
\end{itemize}
The stabiliser of $\vhost + \nhatperp$ is therefore
\begin{equation}\label{eq:Cs_def}
\Cs \;=\; \{e,\,\sigma_E\}, \qquad |\Cs| \;=\; 2.
\end{equation}
This is the elementary mirror group $\Cs$ in Schoenflies notation, the smallest non-trivial point group.

\textbf{Decomposition of $\Reals^4$ at $\vhost$ under $\Cs$.} The reflection $\sigma_E$ acts on $\Reals^4|_{\vhost}$ as the linear hyperplane reflection in the 3-dimensional hyperplane $H^{(3)} = (\ui - \uj)^\perp$. Equivalently, $\sigma_E$ is the identity on $H^{(3)}$ and the sign-flip $-1$ on the 1-dimensional line $\Reals(\ui - \uj)$. The four-dimensional representation of $\Cs$ on $\Reals^4|_{\vhost}$ decomposes as
\begin{equation}\label{eq:Cs_decomp}
\Reals^4 \big|_{\vhost} \;=\; \underbrace{H^{(3)}}_{\Aprime^{\oplus 3}} \;\oplus\; \underbrace{\Reals(\ui - \uj)}_{\Aprpr^{\oplus 1}},
\end{equation}
where $\Aprime$ is the trivial 1-dimensional representation (fixed by $\sigma_E$) and $\Aprpr$ is the non-trivial 1-dimensional representation (flipped by $\sigma_E$). The trivial representation $\Aprime$ has multiplicity~3 in $\Reals^4|_{\vhost}$ under $\Cs$.

\textbf{Three-vector basis for the invariant subspace.} A natural orthogonal basis for the 3-dimensional $\Cs$-invariant subspace $H^{(3)} = (\ui - \uj)^\perp$ at $\vhost$ is $\{\nhatrho,\,\nhatmid,\,\nhatperp\}$:
\begin{itemize}[leftmargin=*]
    \item \emph{$\nhatrho \in H^{(3)}$}: $\nhatrho \cdot (\ui - \uj) = \vhost \cdot \ui - \vhost \cdot \uj$. Since $\ui$ and $\uj$ are both first-shell neighbours of $\vhost$ at the same distance $1/\phi$, $\vhost \cdot \ui = \vhost \cdot \uj$ by symmetry, hence $\nhatrho \perp (\ui - \uj)$. ✓
    \item \emph{$\nhatmid \in H^{(3)}$}: $\nhatmid \propto (\ui + \uj - 2\vhost)$, and $(\ui + \uj) \cdot (\ui - \uj) = |\ui|^2 - |\uj|^2 = 1 - 1 = 0$, plus $\vhost \cdot (\ui - \uj) = 0$ as above. ✓
    \item \emph{$\nhatperp \in H^{(3)}$}: $\nhatperp$ is by construction perpendicular to the face plane, which contains $\ui - \uj$. ✓
    \item \emph{Linear independence}: $\nhatrho$ along $\vhost$ (one specific vertex direction); $\nhatmid$ tangent to $S^3$ at $\vhost$, in the face plane; $\nhatperp$ perpendicular to the face plane. These span three mutually orthogonal directions in $\Reals^4$. ✓
\end{itemize}

\textbf{Conclusion of setup.} The space of $\Cs$-invariant vectors in $\Reals^4$ at $\vhost$ is the 3-dimensional hyperplane
\begin{equation}\label{eq:Sigma_def}
\Sigma_{\Cs} \;:=\; (\ui - \uj)^\perp \;=\; \operatorname{span}\bigl\{\nhatrho,\,\nhatmid,\,\nhatperp\bigr\}.
\end{equation}

\subsection{Main theorem}\label{subsec:main_theorem}

\begin{theorem}[Face-Aligned Invariant-Subspace Structural Theorem]\label{thm:main}
Under hypotheses \emph{(H1$_F$)--(H5$_F$)} as stated in \S\ref{sec:hypotheses}, the net DI-bit current at the host vertex $\vhost$ at every non-zero order $k \geq 1$ in the Mechanism~A perturbation parameter $\delta$ lies in the 3-dimensional $\Cs$-invariant subspace of $\Reals^4$ at $\vhost$:
\begin{equation}\label{eq:main_result}
\jDI(\vhost)\big|_{\Oof{\delta^k}} \;\in\; \Sigma_{\Cs} \;=\; \operatorname{span}\bigl\{\nhatrho,\,\nhatmid,\,\nhatperp\bigr\} \qquad \text{for all } k \geq 1.
\end{equation}
Equivalently, there exist real coefficients $\alpha_k^{(\rho)}$, $\alpha_k^{(\text{ax})}$, $\alpha_k^{(\perp)}$ such that
\begin{equation}\label{eq:three_coeffs}
\jDI(\vhost)\big|_{\Oof{\delta^k}} \;=\; \alpha_k^{(\rho)}\,\nhatrho \;+\; \alpha_k^{(\text{ax})}\,\nhatmid \;+\; \alpha_k^{(\perp)}\,\nhatperp.
\end{equation}
The only symmetry-forced constraint is that the $\ui - \uj$ component of $\jDI(\vhost)|_{\Oof{\delta^k}}$ vanishes identically.
\end{theorem}

This is the weakest of the three Reading-C variants' structural constraints. The symmetry-only argument constrains $\jDI(\vhost)$ to a 3D hyperplane out of $\Reals^4$, leaving 3 of 4 degrees of freedom. Whether the radial coefficient $\alpha_k^{(\rho)}$, the in-face axis coefficient $\alpha_k^{(\text{ax})}$, and/or the face-perpendicular coefficient $\alpha_k^{(\perp)}$ are further forced to vanish (or take specific relations) by Mechanism~A framework-axiom-level constraints is the substantive open question for any future Sequence-2B coefficient enumeration. This is registered as exclusion class~E5 in \S\ref{sec:exclusion}.

\section{Hypotheses}\label{sec:hypotheses}

\begin{description}
\item[(H1$_F$) Face-aligned Reading C:] The substrate-direction primitive $\nhat$ is identified with $\nhatperp$ of \eqref{eq:nperp}, the unit vector at $\vhost$ perpendicular to the chosen triangular face $F = \{\vhost, \ui, \uj\}$, oriented from polytope centre to face centroid.
\item[(H2) 600-cell unit-vertex normalisation:] $|\vhost| = 1$; first-shell neighbours $\ui, \uj$ at distance $1/\phi$ from $\vhost$; standard 600-cell short-edge normalisation.
\item[(H3$_F$) Face-aligned residual symmetry:] The residual symmetry group at $\vhost$ for the face-aligned reading is $\Cs = \{e, \sigma_E\}$ as identified in \S\ref{subsec:stabiliser}. The single non-trivial element $\sigma_E$ is the hyperplane reflection in the 3-hyperplane $(\ui - \uj)^\perp$ through the midpoint of $\{\ui, \uj\}$.
\item[(H4) F.1 v1.0 Theorem~6.1 perturbation-locality:] As in Sequence-1A (Patch~0596 / THEO-DSL-1 / Patch~0550), inherited without modification.
\item[(H5$_F$) Mechanism~A perturbation at face-aligned $\nhat$:] Each oriented edge $\hat{e}$ of the 600-cell edge graph carries the Mechanism~A rate
\begin{equation}\label{eq:H5F}
r(\hat{e}) \;=\; r_0\bigl(1 + \delta\,\hat{e}\cdot\nhatperp\bigr)
\end{equation}
at first order in $\delta$, with $\Oof{\delta^k}$ extension construed via the standard path-integral perturbation-theory machinery per DG-F15-4 default~(A) at sketch level.
\end{description}

The hypothesis structure is parallel to Sequence-1A (Patch~0596 \S{}2), with $\nhatperp$ substituted for $\nhatedge$, $\Cs$ substituted for $\Dfive$, and (H1$_F$, H3$_F$, H5$_F$) substituted for the edge-aligned variants. (H2) and (H4) are unchanged.

\section{Lemma~1: $\Cs$-equivariance of the substrate current at $\vhost$}\label{sec:lemma1}

\begin{lemma}[Order-by-order $\Cs$-invariance]\label{lem:cs_invariance}
Under (H1$_F$)--(H5$_F$), for every $k \geq 0$,
\begin{equation}\label{eq:cs_inv}
\sigma_E \cdot \jDI(\vhost)\big|_{\Oof{\delta^k}} \;=\; \jDI(\vhost)\big|_{\Oof{\delta^k}}.
\end{equation}
\end{lemma}

\begin{proof}
The argument follows the three-step pattern of Patch~0596 Lemma~1 (which used $\Dfive$-equivariance) with $\Cs$ substituted for $\Dfive$. The proof is shorter because $\Cs$ has only one non-trivial element $\sigma_E$ to verify.

\textbf{Step 1 ($\sigma_E$-equivariance of the rate function).} Mechanism~A's rate function on each oriented edge $\hat{e}$ is $r(\hat{e}) = r_0(1 + \delta\,\hat{e}\cdot\nhatperp)$ per (H5$_F$) / \eqref{eq:H5F}. The reflection $\sigma_E$ preserves $\nhatperp$ direction by construction (\S\ref{subsec:stabiliser}). For any oriented edge $\hat{e}$ of the 600-cell edge graph, $\sigma_E\hat{e}\cdot\nhatperp = \hat{e}\cdot\sigma_E^{-1}\nhatperp = \hat{e}\cdot\nhatperp$ (using $O(4)$-invariance of the inner product and $\sigma_E^{-1}\nhatperp = \nhatperp$). Therefore
\begin{equation}\label{eq:rate_eq}
r(\sigma_E\hat{e}) \;=\; r_0\bigl(1 + \delta\,\sigma_E\hat{e}\cdot\nhatperp\bigr) \;=\; r_0\bigl(1 + \delta\,\hat{e}\cdot\nhatperp\bigr) \;=\; r(\hat{e}).
\end{equation}

\textbf{Step 2 ($\sigma_E$-permutation of the path set).} $\sigma_E \in \Cs \subset \Icoh = \text{stab}_{\Hfour}(\vhost)$, so $\sigma_E$ acts on the 600-cell edge graph $\Gsixcell$ as a graph automorphism fixing $\vhost$. By (H4) shell-locality, $\sigma_E \cdot \Ek(\vhost) = \Ek(\vhost)$ as edge sets, hence $\sigma_E$ permutes the path set $\mathcal{P}_k(\vhost)$ as a finite set:
\begin{equation}\label{eq:path_perm}
\sigma_E: \mathcal{P}_k(\vhost) \to \mathcal{P}_k(\vhost), \qquad P \mapsto \sigma_E P.
\end{equation}

\textbf{Step 3 ($O(4)$-equivariance of edge-direction vectors).} For each path $P \in \mathcal{P}_k(\vhost)$, the directed edge-displacement contribution $\vec{d}(P) \in \Reals^4$ satisfies $\vec{d}(\sigma_E P) = \sigma_E \cdot \vec{d}(P)$ by the orthogonal action of $\sigma_E \subset O(4)$ on $\Reals^4$ at $\vhost$.

\textbf{Combining the three steps.} As in Patch~0596 Lemma~1, the substrate current $\Oof{\delta^k}$ sum over paths is re-indexed by the $\sigma_E$ action, giving $\sigma_E \cdot \jDI(\vhost)|_{\Oof{\delta^k}} = \jDI(\vhost)|_{\Oof{\delta^k}}$, which is \eqref{eq:cs_inv}.
\end{proof}

\textbf{Remark on robustness.} The proof uses only $\Cs$-equivariance of Mechanism~A's rate function and $\Cs$-permutation of the path set. As with Patch~0596, the argument is robust across DG-F15-4 sub-options~A (default path-integral), B (alternative framework-axiom MA.3), and C (Layer~4 derivation), since only $\Cs$-equivariance of the path-integral construction is required.

\section{Lemma~2: The $\Cs$-invariant subspace of $\Reals^4$ at $\vhost$}\label{sec:lemma2}

\begin{lemma}[$\Cs$-invariant subspace identification]\label{lem:cs_invariant_subspace}
Under (H1$_F$)--(H3$_F$), the subspace of $\Cs$-invariant vectors in $\Reals^4$ at $\vhost$ is the 3-dimensional hyperplane
\begin{equation}\label{eq:Sigma_value}
\Sigma_{\Cs} \;=\; (\ui - \uj)^\perp \;=\; \operatorname{span}\bigl\{\nhatrho,\,\nhatmid,\,\nhatperp\bigr\}.
\end{equation}
\end{lemma}

\begin{proof}
The argument is direct from the decomposition \eqref{eq:Cs_decomp}: $\Reals^4|_{\vhost}$ decomposes under $\Cs$ as $\Aprime^{\oplus 3} \oplus \Aprpr^{\oplus 1}$, where the trivial representation $\Aprime$ has multiplicity~3 occupying the hyperplane $(\ui - \uj)^\perp$ and the non-trivial representation $\Aprpr$ has multiplicity~1 occupying the line $\Reals(\ui - \uj)$. The space of $\Cs$-invariant vectors is exactly the multiplicity space of $\Aprime$, which is the 3-dimensional hyperplane $(\ui - \uj)^\perp$. The basis $\{\nhatrho, \nhatmid, \nhatperp\}$ for this hyperplane was verified at \S\ref{subsec:stabiliser}: each of these three directions is in $(\ui - \uj)^\perp$ (verified inner products) and they are mutually linearly independent (radial vs.\ in-face-axis vs.\ face-perpendicular).

Equivalently via character calculation: $\Cs$ has two irreducible representations $\Aprime$ (1-dimensional trivial, $\chi(e) = 1, \chi(\sigma_E) = 1$) and $\Aprpr$ (1-dimensional sign, $\chi(e) = 1, \chi(\sigma_E) = -1$). On $\Reals^4|_{\vhost}$, the character is $\chi(e) = 4$, $\chi(\sigma_E) = \text{trace}(\sigma_E) = 3 - 1 = 2$ (three $+1$ eigenvalues from the 3-hyperplane fixed by reflection plus one $-1$ from the reflected line). The multiplicity of $\Aprime$ is $m_{\Aprime} = (1\cdot 4 + 1\cdot 2)/2 = 3$. The multiplicity of $\Aprpr$ is $m_{\Aprpr} = (1\cdot 4 + (-1)\cdot 2)/2 = 1$. Sum $3 + 1 = 4 = \dim \Reals^4$. ✓
\end{proof}

\section{Main proof of Theorem~\ref{thm:main}}\label{sec:main_proof}

\begin{proof}[Proof of Theorem~\ref{thm:main}]
By Lemma~\ref{lem:cs_invariance}, $\jDI(\vhost)|_{\Oof{\delta^k}}$ is $\Cs$-invariant for every $k \geq 1$. By Lemma~\ref{lem:cs_invariant_subspace}, the $\Cs$-invariant subspace of $\Reals^4|_{\vhost}$ is $\Sigma_{\Cs} = \operatorname{span}\{\nhatrho, \nhatmid, \nhatperp\}$. Therefore $\jDI(\vhost)|_{\Oof{\delta^k}} \in \Sigma_{\Cs}$, which is \eqref{eq:main_result}. The decomposition \eqref{eq:three_coeffs} follows from the 3-dimensionality of $\Sigma_{\Cs}$ and the chosen basis.
\end{proof}

\begin{corollary}[All-orders invariant-subspace strengthening]\label{cor:all_orders}
As in Sequence-1A, the hypotheses (H1$_F$)--(H5$_F$) make no use of $k = 2$ specifically; the invariant-subspace conclusion holds at every $k \geq 1$. The full formal $\delta$-power series at $\vhost$ is contained in $\Sigma_{\Cs}$ order-by-order:
\begin{equation}
\jDI(\vhost; \delta) \;=\; \sum_{k \geq 1} \delta^k \bigl(\alpha_k^{(\rho)}\,\nhatrho + \alpha_k^{(\text{ax})}\,\nhatmid + \alpha_k^{(\perp)}\,\nhatperp\bigr).
\end{equation}
Any future Sequence-2B coefficient enumeration under (H5$_F$) may impose this 3-vector ansatz from the outset, reducing the computational target from a 4-vector to three scalars $\alpha_k^{(\rho)}, \alpha_k^{(\text{ax})}, \alpha_k^{(\perp)}$ at each order $k$.
\end{corollary}

\section{Three-variant comparison table and the dimensional-growth pattern}\label{sec:comparison}

The completed Sequence-1 work across both non-vertex-aligned variants now allows the full three-variant comparison:

\begin{center}
\begin{tabular}{lcll}
\textbf{Reading C variant} & \textbf{Stabiliser} & \textbf{Order} & \textbf{Invariant subspace of $\Reals^4|_{\vhost}$} \\
\hline
Vertex-aligned (THEO-DSL-4 / Patch 0585) & $\Icoh$ & 120 & $\Reals\nhatrho$ (1D) \\
Edge-aligned (THEO-DSL-6 / Patch 0596) & $\Dfive$ & 10 & $\operatorname{span}\{\nhatrho, \nhatedge\}$ (2D) \\
Face-aligned (THEO-DSL-8 / Patch 0597; this artifact) & $\Cs$ & 2 & $\operatorname{span}\{\nhatrho, \nhatmid, \nhatperp\}$ (3D) \\
\end{tabular}
\end{center}

\textbf{The dimensional-growth pattern.} The invariant subspace dimension grows monotonically as the stabiliser shrinks: $1 \to 2 \to 3$ as the stabiliser order falls $120 \to 10 \to 2$. The pattern is consistent with the group-theoretic principle that the multiplicity of the trivial representation in the standard 4D representation of $G \subset O(4)$ grows as $G$ shrinks; in the extreme case $G = \{e\}$ all four dimensions would be ``$\Aprime$-trivial'' and the invariant subspace would be the entire $\Reals^4$.

\textbf{Vertex-aligned uniqueness.} The 1-dimensional invariant subspace at vertex-aligned Reading C is the only one of the three that produces a strict parallel-to-$\nhat$ structure for the substrate current $\jDI(\vhost)$. The 2-dimensional edge-aligned invariant subspace admits a radial-component degree of freedom; the 3-dimensional face-aligned invariant subspace admits two additional in-face degrees of freedom. The substrate-locality umbrella theorem of F.1 (Theorem~7.1) is therefore a vertex-aligned-Reading-C-specific consequence of the maximal residual symmetry $\Icoh$. \textbf{This converges with Capotauro v2.0's vertex-aligned-uniqueness selection at Q1$'$+Q1$'$.A (Patch~0419 Finding C-W37) at the F.1 temporal-sector level}: under non-vertex-aligned Reading-C variants, the substrate-locality structure substantively weakens.

\textbf{Cross-sector consistency.} Both Capotauro v2.0 (spatial-sector chirality) and F.1 (temporal-sector substrate-locality) now have established structural uniqueness arguments for vertex-aligned Reading C: Capotauro at the cage-shell averaging factor level ($\chi/6$ specific to $\Icoh$-equivariant cage-shell sums) and F.1 at the parallel-to-$\nhat$ substrate-current level. The two arguments converge on the same conclusion through structurally distinct routes, strengthening the Reading-C selection.

\section{Anti-erasure remarks}\label{sec:antierasure}

\subsection{Anti-erasure: stabiliser at face-aligned Reading C relative to FP.md}\label{subsec:antierasure_FPmd}

\begin{remark}[Anti-erasure: stabiliser at face-aligned Reading C]\label{rem:antierasure_stabiliser}
The frontier-sector entry \texttt{frontier\_sectors/FP.md} line~230 states, for the face-aligned Reading-C variant: ``residual $D_2$ symmetry at face centroid.'' This identification is incorrect under the vertex-host convention (preserved from Sequence-1A precedent, Patch~0596) and is registered openly per the CPP programme's anti-erasure discipline (cf.\ Patches~0596 Remark~5.1 + 5.2; precedent Patches 0587 + 0590 + 0591):

\begin{enumerate}[label=(\arabic*)]
    \item \textbf{Incorrect group identification under vertex-host convention.} Under the vertex-host convention (the natural F.1-paper-consistent choice), the stabiliser of $\vhost$~+~$\nhatperp$ direction is $\Cs$ of order~2, not $D_2$ of order~4. The order-2 stabiliser consists of the identity plus the single edge-swap reflection $\sigma_E$ in the 3-hyperplane $(\ui - \uj)^\perp$.
    \item \textbf{Alternative interpretation under face-centroid-host convention.} If instead the host were taken to be the face centroid (not a 600-cell vertex), the stabiliser of the centroid~+~$\nhatperp$ would be the full $D_3$ of order~6 (the triangle symmetry group acting on the 3 face vertices, with the face-perp direction preserved). However, this alternative interpretation breaks the F.1-paper-consistent vertex-host convention and introduces cross-host complications (different graph structure, different shell decomposition, modified (H4) graph-extension); it was rejected at Patch~0595 \S{}1.2 and at \S\ref{subsec:setup} of the present Patch.
    \item \textbf{Neither identification matches FP.md's ``$D_2$''.} Neither $\Cs$ (vertex-host) nor $D_3$ (face-centroid-host) matches the FP.md ``$D_2$'' label. The error is registered here without silent correction; the follow-on registry-propagation Patch (candidate 0597a) will update FP.md line~230 to reflect $\Cs$ at face-aligned under the vertex-host convention.
\end{enumerate}

The substantive structural conclusion of Theorem~\ref{thm:main} (3-dimensional $\Cs$-invariant subspace = $(\ui - \uj)^\perp$) is specific to the vertex-host convention and the corrected $\Cs$ stabiliser identification.
\end{remark}

\subsection{Anti-erasure: vertex-host convention consistency}\label{subsec:vertex_host_consistency}

\begin{remark}[Vertex-host convention preserved across all three Reading-C variants]\label{rem:vertex_host}
The three Reading-C variant artifacts (Patches~0585, 0596, 0597) all preserve the vertex-host convention with $\vhost$ as a 600-cell vertex; the variation is in the choice of $\nhat$ direction at the host. This is the natural F.1-paper-consistent setup and avoids cross-host complications. The alternative face-centroid-host setup (or edge-midpoint-host setup) is a separate research question not addressed by the present OPEN-FP-F1-5 trajectory; if pursued, it would constitute a new Open Problem (potentially OPEN-FP-F1-7 or similar) registered separately. The present Patch closes OPEN-FP-F1-5 Sequence-1B under the vertex-host convention only.
\end{remark}

\section{Five-class exclusion enumeration}\label{sec:exclusion}

\paragraph{(E1) Sequence-2B coefficient sub-question NOT closed.} Theorem~\ref{thm:main} is symmetry-only and does not compute the closed-form coefficients $\alpha_k^{(\rho)}, \alpha_k^{(\text{ax})}, \alpha_k^{(\perp)}$. Coefficient computation under (H5$_F$) is the target of Sequence-2B (Patches~0598+).

\paragraph{(E2) Edge-aligned variant separately closed.} The edge-aligned Reading-C analog is the target of Sequence-1A (Patch~0596 / THEO-DSL-6) and is now CLOSED at publication-grade Layer~3 unconditional. Both Sequence-1 structural findings (edge-aligned 2D, face-aligned 3D) are now in hand, completing DG-F15-1 default-(C) priority.

\paragraph{(E3) (H5$_F$) framework-extension at sketch-document Layer~3.} As with (H5) and (H5$_E$), the Mechanism~A perturbation at $\Oof{\delta^k}$ for $k \geq 2$ under (H5$_F$) is at sketch-document Layer~3 per DG-F15-4 default~(A). Publication-grade hardening of (H5$_F$) is a follow-on question parallel to the F1-1 priority queue item~(G).

\paragraph{(E4) Layer~4 axiomatic derivation NOT in scope.} Inherits the shared scope-line of the F.1 Layer~3 hardened-theorem sequence: OPEN-FP-F1-2 is the long-term programme target. MA.1 + MA.2 are taken as framework axioms.

\paragraph{(E5) Three-coefficient question open at framework-axiom level.} Theorem~\ref{thm:main}'s 3-dimensional invariant subspace admits three substantive coefficients $\alpha_k^{(\rho)}, \alpha_k^{(\text{ax})}, \alpha_k^{(\perp)}$ on symmetry grounds alone. Whether any of these are independently forced to vanish, or whether they satisfy specific relations beyond symmetry (e.g.\ via cage-shell-averaging considerations at the framework level, or via explicit path-class enumeration under (H5$_F$)), is NOT addressed by the present artifact and is the natural substantive target of Sequence-2B's coefficient enumeration. The substantive richness of the open question is greater for face-aligned than edge-aligned (three coefficients vs.\ two; cf.\ Patch~0596 E5).

\section{Cross-trajectory linkage and registration}\label{sec:registration}

\textbf{Cross-trajectory linkage.} This artifact is the \textbf{thirteenth} in the F.1 Layer~3 hardened-theorem sequence. Combined approximate length across the thirteen \texttt{.tex} artifacts is now $\sim$4{,}230 lines LaTeX. The OPEN-FP-F1-5 trajectory now has two structural theorems closed at publication-grade Layer~3 unconditional (THEO-DSL-6 + THEO-DSL-8 candidates); the four-artifact trajectory completion estimated at Patch~0595 \S{}1.2 is approximately at its midpoint (2 of 4 Patches landed; 2 Sequence-2 coefficient Patches remaining).

\textbf{Theorem-registry naming.} Per DG-F15-5 default~(A), the present result is registered as candidate \textbf{THEO-DSL-8} entry. (THEO-DSL-7 is reserved for the Sequence-2A candidate-coefficient closure at the edge-aligned variant; the registry slot ordering preserves the (structural, coefficient) pairing convention: THEO-DSL-6 + THEO-DSL-7 for edge-aligned, THEO-DSL-8 + THEO-DSL-9 for face-aligned.) Registry propagation to \texttt{theorem-registry.md} + \texttt{frontier\_sectors/FP.md} (including the stabiliser-identification correction per Remark~\ref{rem:antierasure_stabiliser}) + \texttt{master\_glossary.md} (new term entries for $\nhatperp$, $\nhatmid$, $\Cs$, face-aligned Reading C, Sequence-1B) is deferred to a follow-on registry-propagation Patch (candidate 0597a).

\textbf{Cross-sector cross-check.} The face-aligned variant's 3-dimensional invariant subspace represents the weakest structural constraint of the three variants. Combined with Capotauro v2.0's vertex-aligned-uniqueness selection, the face-aligned variant emerges as structurally inadequate for substrate-locality in both the spatial sector (Capotauro chirality-matrix-element selection) and the temporal sector (F.1 substrate-current direction). This cross-sector convergence on vertex-aligned uniqueness through structurally distinct routes is the central physical result of the present Patch.

\section{Falsifier and honest scope-limitation framing}\label{sec:falsifier}

\textbf{Falsifier.} Demonstration of $\jDI(\vhost)|_{\Oof{\delta^k}}$ having a component along $\ui - \uj$ at any order $k \geq 1$, from a valid first-principles $\Cs$-equivariant calculation under (H1$_F$)--(H5$_F$), would contradict Lemma~\ref{lem:cs_invariant_subspace}'s identification of the $\Cs$-invariant subspace as $(\ui - \uj)^\perp$. Equivalently, demonstration that Mechanism~A's $\Cs$-equivariance under (H1$_F$)~+~(H3$_F$) fails at any order would contradict Lemma~\ref{lem:cs_invariance}.

\textbf{Honest scope-limitation framing.} Theorem~\ref{thm:main} closes the structural component of OPEN-FP-F1-5 Sequence-1B only; it constrains $\jDI(\vhost)|_{\Oof{\delta^k}}$ to a 3D hyperplane without determining the relative weighting of $\nhatrho, \nhatmid, \nhatperp$ components. The coefficient sub-question (E1) is the target of Sequence-2B; the three-coefficient relative-magnitude question (E5) is open at the framework-axiom level. The (H5$_F$) framework-extension qualifier applies only to follow-on Sequence-2B coefficient work; the present Sequence-1B result is publication-grade Layer~3 unconditional under the vertex-host convention (Remark~\ref{rem:vertex_host}).

\section{References}\label{sec:references}

\textbf{Internal references.} Inherits the reference set of Patch~0596 \S{}9 (F.1 v1.0 paper; THEO-DSL-1 through THEO-DSL-6 candidate; eleven prior hardened-theorem artifacts; scoping document Patch~0595; FP.md OPEN-FP-F1-5 entry). The present artifact adds:
\begin{itemize}[leftmargin=*]
    \item Patch~0596 / \texttt{edge\_aligned\_invariant\_subspace\_structural.tex}: Sequence-1A predecessor establishing THEO-DSL-6 candidate (edge-aligned 2D invariant subspace under $\Dfive$).
\end{itemize}

\textbf{Cross-paper context.} Capotauro v2.0 \S{}2.3 + Patch 0419 Finding C-W37 (Reading-C three-variant framing; vertex-aligned selection). The present Patch's face-aligned 3D invariant subspace result strengthens Capotauro's vertex-aligned-uniqueness selection at the F.1 substrate-locality level.

\textbf{Polytope theory.} Coxeter, H.\ S.\ M.\ \emph{Regular Polytopes} (3rd ed., Dover, 1973): standard 600-cell face structure (1200~triangular faces; 30 faces per vertex via orbit-stabiliser on $\Icoh$).

% BEGIN GENERATED IDENTIFIER APPENDIX -- do not edit by hand
\clearpage
\section*{Appendix: Programme identifiers used in this paper}
\addcontentsline{toc}{section}{Appendix: Programme identifiers}

\noindent This programme tracks its results, assumptions and unresolved questions under short identifiers, so that a claim and its dependencies can be cited precisely. Prefixes denote the kind of item: \texttt{THEO-} a proved theorem, \texttt{FI-} a foundational input the derivation assumes, \texttt{PRED-} a registered prediction, \texttt{CONJ-} a conjecture, \texttt{OPEN-} an unresolved question, \texttt{PH-} a problem history. Those appearing in this paper are glossed below; no external document is needed to read them.

\begin{description}
  \item[\texttt{OPEN-FP-F1-1}] \hfill \\ Extend the F.1 substrate-locality umbrella theorem (Theorem 7.1) to \$\textbackslash\{\}mathcal\{O\}(\textbackslash\{\}delta\textasciicircum{}2)\$ by deriving the second-shell inner-product and edge-projection identities for the 600-cell, and determine whether the parallel-to-\$\textbackslash\{\}hat\{n\}\$ structure of \$\textbackslash\{\}vec\{J\}\_\{\textbackslash\{\}text\{DI-net\}\}(v\_\{\textbackslash\{\}text\{host\}\})\$ survives at second order. **Resolution**: both sub-questions closed — structural sub-question at publication-grade Layer 3 unconditional (\$\textbackslash\{\}vec\{j\}\_k(v\_\{\textbackslash\{\}text\{host\}\}) \textbackslash\{\}parallel \textbackslash\{\}hat\{n\}\$ at every \$k \textbackslash\{\}geq 1\$ via THEO-DSL-4) + coefficient sub-question at sketch-document Layer 3 under (H5) (\$\textbackslash\{\}alpha\_2 = -9/\textbackslash\{\}phi\textasciicircum{}2 = -9(2-\textbackslash\{\}phi) \textbackslash\{\}approx -3.438\$ via THEO-DSL-5 candidate; ratio \$\textbackslash\{\}alpha\_2/\textbackslash\{\}alpha\_1 = -3/2\$ exactly).
  \item[\texttt{OPEN-FP-F1-2}] \hfill \\ Derive Mechanism A (F.1 framework axioms MA.1 + MA.2: propagation-rate asymmetry \$r(\textbackslash\{\}hat\{e\}) = r\_0(1 + \textbackslash\{\}delta\textbackslash\{\},\textbackslash\{\}hat\{e\}\textbackslash\{\}cdot\textbackslash\{\}hat\{n\})\$ and framework-local current construction at \$\textbackslash\{\}mathcal\{O\}(\textbackslash\{\}delta\textasciicircum{}1)\$) from CPP primitive axioms A1–A11 alone.
  \item[\texttt{OPEN-FP-F1-5}] \hfill \\ Extend or reformulate F.1 Theorems 5.1 (host-first-shell projection), 5.2 (first-shell perpendicularity), and 7.1 (substrate-locality umbrella) under edge-aligned Reading C (\$\textbackslash\{\}hat\{n\}\$ along a 600-cell short-edge direction; residual \$D\_5\$ symmetry at edge midpoint per Patch 0596 Remark 5.1 anti-erasure correction — see below) and face-aligned Reading C (\$\textbackslash\{\}hat\{n\}\$ perpendicular to a triangular face; residual \$C\_s\$ symmetry under vertex-host convention per Patch 0597 Remark 7.1 anti-erasure correction — see below).
  \item[\texttt{OPEN-FP-F1-7}] \hfill \\ A reserved identifier for a possible new open problem beyond the F1-5 trajectory, should higher-order extensions be pursued.
  \item[\texttt{OPEN-SD-CHIR-PRIMITIVE}] \hfill \\ Derive the universe's primitive chirality bias from a single substrate-level mechanism (current leading candidate: primitive 4D direction \$\textbackslash\{\}hat\{n\}\$ aligned with a 600-cell host vertex per Reading C) that simultaneously accounts for the five observable manifestations across the CPP corpus: K3-doublet chirality (Capotauro \$\textbackslash\{\}Delta p\_\{LR\}\$, mass-mixing sector), W-bracelet V-A coupling (electroweak, massless helicity limit), electromagnetic handedness (E×B / Poynting / Maxwell right-hand rule and qDP/eDP polarization patterns), thermodynamic causal arrow (DI-bit propagation direction / PCD cycle ordering / Conscious Moment sequence), and cosmological vacuum asymmetry (matter-antimatter / baryogenesis / Capotauro nucleation-event sign-selection).
  \item[\texttt{THEO-DSL-1}] \hfill \\ Perturbation-Theory Propagation Rule + Shell-Locality at \$\textbackslash\{\}mathcal\{O\}(\textbackslash\{\}delta\textasciicircum{}1)\$ (**first F-line flagship theorem under THEO-DSL-N sub-prefix convention**; publication-grade Layer 3 **unconditional**; first of three F.1 v1.0 SHIPPED theorems closing OPEN-SD-CHIR-PRIMITIVE manifestation (iv) thermodynamic causal arrow at sketch-document Layer 3 via Theorem 7.1 substrate-locality umbrella)
  \item[\texttt{THEO-DSL-4}] \hfill \\ \$\textbackslash\{\}mathcal\{O\}(\textbackslash\{\}delta\textasciicircum{}k)\$ Substrate-Current Parallel-to-\$\textbackslash\{\}hat\{n\}\$ Structural Theorem at All Orders \$k \textbackslash\{\}geq 1\$ (**fourth F-line flagship theorem under THEO-DSL-N sub-prefix convention**; publication-grade Layer 3 **unconditional** for the structural sub-question; **first Route C sequence 1 artifact** under the OPEN-FP-F1-1 trajectory; **closes the structural sub-question of OPEN-FP-F1-1** at all orders \$k \textbackslash\{\}geq 1\$ via symmetry-only argument; coefficient \$\textbackslash\{\}alpha\_k\$ at \$k \textbackslash\{\}geq 2\$ remains OPEN as target of Route C sequence 2)
  \item[\texttt{THEO-DSL-6}] \hfill \\ The order-10 stabilizer of an oriented first-shell edge in the icosahedral group at the host vertex.
  \item[\texttt{THEO-DSL-7}] \hfill \\ Umbrella registration covering first- and second-order edge-aligned coefficient content in the Dynamical Substrate Law.
  \item[\texttt{THEO-DSL-8}] \hfill \\ The cross-shell edge orbit closure under the reflection stabilizer of face-aligned Reading C.
  \item[\texttt{THEO-DSL-9}] \hfill \\ The candidate coefficient theorem for the face-aligned branch of the substrate law.
  \item[\texttt{THEO-DSL-N}] \hfill \\ (family) A proposed family registration for the shell-perpendicularity pattern that recurs at every perturbative order.
\end{description}

\subsection*{How we review}

\noindent Results in this programme are checked by an \emph{AI review panel}: several large language models from different vendors are given the same paper and the same fixed protocol, and asked to find errors and raise objections independently. Objections are recorded and either answered in the text or registered as open problems under the identifiers glossed above.

\noindent This is a \emph{verification method the authors apply to their own work}. It is not journal peer review, it is not equivalent to it, and agreement among panel members is not evidence that a result is correct --- only that no member found a specific fault. Where individual models are named in the text, they are named for reproducibility of the method; model versions change, and a reader should treat the named models as a record of what was run rather than as an endorsement.

% END GENERATED IDENTIFIER APPENDIX

\end{document}
